%! AP Statistics — Unit 4, Lesson 4.9 — Homework (Answer Key)
\documentclass[10pt]{article}
\usepackage{apstats-article}
\usepackage{apstats-key}
\usepackage{grading}
\gradingassignment{apstats-u04-l09-hw}
\gradingcourse{advanced_statistics}
\newcommand{\TallMath}[1]{$\displaystyle #1\rule[-1.4em]{0pt}{3.2em}$}

\begin{document}

\pageheader{Unit 4, Lesson 4.9}{Homework --- Answer Key}
\namedateperiod

\vspace{0.1in}

\begin{practicebox}
\small

\textbf{1.} $W = 3 + 5X$; $\mu_X = 20$, $\sigma_X = 4$.

\begin{enumerate}[label=(\alph*), leftmargin=2em, itemsep=8pt]
  \item $\mu_W = $ \ans{$3 + 5(20) = 103$}\gradedkey{\smash{\raisebox{-4pt}{\gradedblank[id=hw1a, type=short, points=1, answer=103, check=numeric, tolerance=0]{2.4cm}{0.2in}}}}
  \item $\sigma_W = $ \ans{$|5|(4) = 20$}\gradedkey{\smash{\raisebox{-4pt}{\gradedblank[id=hw1b, type=short, points=1, answer=20, check=numeric, tolerance=0]{2.4cm}{0.2in}}}}
  \item \ansline{$\mu_W$ increases by $5 \times 2 = 10$ (to 113).}\gradedkey{\smash{\raisebox{-4pt}{\gradedblank[id=hw1c, type=short, points=1, answer=10, check=numeric, tolerance=0, display={increases by 10}]{3.4cm}{0.2in}}}}
\end{enumerate}

\vspace{0.14in}

\textbf{2.} Pizza and soda.

\begin{enumerate}[label=(\alph*), leftmargin=2em, itemsep=8pt]
  \item \ans{$\mu_{X+Y} = \$12+\$3 = \$15.00$.} \ansline{On average, a customer spends \$15.00 on a pizza and soda combined.}\gradedkeyline{$\mu_{X+Y} = $~\gradedblank[id=hw2a1, type=short, points=1, answer=15, check=numeric, tolerance=0.01, display={\$15.00}]{0.9in}{0.2in}\hfill \gradedblank[id=hw2a2, type=rubric, points=1, display={On average, a customer spends about \$15.00 on a pizza and soda combined.}]{4.9in}{\gradedlinesheight{2}-5pt}}

  \item \ans{$\sigma^2_{X+Y}=2.00^2+0.75^2=4.00+0.5625=4.5625$; $\sigma_{X+Y}=\sqrt{4.5625}\approx\$2.14$}\gradedkey{\raisebox{-3pt}{$\sigma_{X+Y} = $~\gradedblank[id=hw2b, type=short, points=2, answer=2.14, check=numeric, tolerance=0.01]{0.9in}{0.2in}}}

  \item \ansline{The friend added standard deviations directly (\$2.00+\$0.75=\$2.75); you must add variances first, then take the square root.}\gradedkey{\raisebox{-3pt}{\gradedblank[id=hw2c, type=rubric, points=2, display={SDs don't add: add variances, then take the square root.}]{4.4in}{0.22in}}}
\end{enumerate}

\vspace{0.14in}

\textbf{3.} Bowling.

\begin{enumerate}[label=(\alph*), leftmargin=2em, itemsep=8pt]
  \item \ans{$\mu_{A-B}=140-120=20$; $\sigma^2_{A-B}=18^2+15^2=324+225=549$; $\sigma_{A-B}=\sqrt{549}\approx 23.4$}\hfill\gradedright{\smash{\raisebox{18pt}{$\mu_{A-B} = $~\gradedblank[id=hw3a1, type=short, points=1, answer=20, check=numeric, tolerance=0]{0.9in}{0.2in}\quad $\sigma_{A-B} = $~\gradedblank[id=hw3a2, type=short, points=2, answer=23.4, check=numeric, tolerance=0.1]{0.9in}{0.2in}}}}

  \item \ansline{On average, player A scores about 20 pins more than player B per game.}\gradedkeyline{\kern0.4in\gradedblank[id=hw3b, type=rubric, points=2, display={On average, player A scores about 20 pins more than player B per game.}]{\linewidth+6pt-0.4in}{\gradedlinesheight{2}-5pt}}

  \item \ansline{When subtracting two independent RVs, both contribute uncertainty; variances still add, giving a larger SD than either player's individual SD.}\gradedkey{\raisebox{-8pt}{\gradedblank[id=hw3c, type=rubric, points=2, display={Variances add even for a difference: both add variability, so $\sigma_{A-B} = 23.4$.}]{3.2in}{\gradedlinesheight{2}-5pt}}}
\end{enumerate}

\vspace{0.14in}

\textbf{4.} $\mu_X=5$, $\sigma_X=2$, $\mu_Y=3$, $\sigma_Y=4$.

\begin{enumerate}[label=(\alph*), leftmargin=2em, itemsep=8pt]
  \item $\mu_{2X-Y} = $ \ans{$2(5)-3=7$}~\smash{\raisebox{-4pt}{\gradedblank[id=hw4a, type=short, points=1, answer=7, check=numeric, tolerance=0]{2.4cm}{0.2in}}}
  \item \ans{$\sigma^2_{2X-Y}=(2)^2\sigma^2_X+\sigma^2_Y=4(4)+16=16+16=32$; $\sigma_{2X-Y}=\sqrt{32}=4\sqrt{2}\approx 5.66$}\hfill\gradedright{\smash{\raisebox{20pt}{\gradedblank[id=hw4b, type=short, points=2, answer=5.66, check=numeric, tolerance=0.01]{2.3cm}{0.2in}}}}
\end{enumerate}

\end{practicebox}

\vspace{0.1in}

\begin{extensionbox}
\small
\textbf{Extension.}
\begin{enumerate}[leftmargin=1.5em, itemsep=6pt]
  \item \ans{$\mu_T=100+80=180$; $\sigma^2_T=15^2+12^2=225+144=369$; $\sigma_T=\sqrt{369}\approx 19.21$. $T \sim N(180,19.21)$.}

  \item \ans{$z=(200-180)/19.21\approx 1.04$; $P(T>200)=P(Z>1.04)\approx 0.1492$.}
\end{enumerate}
\end{extensionbox}

\vspace{0.1in}

\begin{blurbbox}[Preview: Lesson 4.10]{sky}
\small
In Lesson 4.10 you will meet the \textbf{binomial distribution} --- a special discrete
distribution for the count of ``successes'' in $n$ independent trials. Its mean $\mu = np$
and variance $\sigma^2 = np(1-p)$ follow directly from the rules you learned today.
\end{blurbbox}

\end{document}
